\begin{tabular}{llp{3in}}

	$\mathbf{A}$           & \hspace{1.5cm} = & matriks sistem (\textit{state matrix}) \\
	$D$                    & \hspace{1.5cm} = & koefisien redaman (p.u.) \\
	$f_d$                  & \hspace{1.5cm} = & frekuensi osilasi teredam (Hz) \\
	$H$                    & \hspace{1.5cm} = & konstanta inersia (MWs/MVA) \\
	$K_s$                  & \hspace{1.5cm} = & koefisien daya sinkronisasi (p.u./rad) \\
	$P_e$                  & \hspace{1.5cm} = & daya elektrik keluaran (p.u.) \\
	$P_m$                  & \hspace{1.5cm} = & daya mekanik masukan (p.u.) \\
	$P_{ki}$               & \hspace{1.5cm} = & \textit{participation factor} variabel keadaan ke-$k$ pada mode ke-$i$ \\
	$\delta$               & \hspace{1.5cm} = & sudut rotor generator sinkron (rad) \\
	$\Delta\delta_{12}$    & \hspace{1.5cm} = & selisih sudut poros turbin--generator DFIG-WT (rad) \\
	$\zeta$                & \hspace{1.5cm} = & \textit{damping ratio} \\
	$\lambda$              & \hspace{1.5cm} = & \textit{eigenvalue} \\
	$\boldsymbol{\Lambda}$ & \hspace{1.5cm} = & matriks diagonal \textit{eigenvalue} \\
	$\phi_{ki}$            & \hspace{1.5cm} = & elemen ke-$k$ dari \textit{right eigenvector} mode ke-$i$ \\
	$\boldsymbol{\Phi}$    & \hspace{1.5cm} = & matriks \textit{right eigenvector} \\
	$\varphi_m$            & \hspace{1.5cm} = & sudut rotor internal DFIG-WT (rad) \\
	$\psi_{ik}$            & \hspace{1.5cm} = & elemen ke-$k$ dari \textit{left eigenvector} mode ke-$i$ \\
	$\boldsymbol{\Psi}$    & \hspace{1.5cm} = & matriks \textit{left eigenvector} \\
	$\psi_{fd}$            & \hspace{1.5cm} = & fluks kumparan medan eksitasi generator sinkron \\
	$\psi_{1d}$            & \hspace{1.5cm} = & fluks kumparan peredam sumbu-$d$ generator sinkron \\
	$\psi_{1q},\,\psi_{2q}$ & \hspace{1.5cm} = & fluks kumparan peredam sumbu-$q$ generator sinkron \\
	$\sigma$               & \hspace{1.5cm} = & bagian riil \textit{eigenvalue} \\
	$\omega$               & \hspace{1.5cm} = & kecepatan sudut rotor generator sinkron (rad/s) \\
	$\omega_d$             & \hspace{1.5cm} = & frekuensi sudut osilasi teredam (rad/s) \\
	$\omega_r$             & \hspace{1.5cm} = & kecepatan sudut rotor DFIG-WT (rad/s) \\
	$\omega_s$             & \hspace{1.5cm} = & kecepatan sudut sinkron (rad/s) \\
\end{tabular}

\newpage

\begin{tabular}{llp{3in}}

	AVR      & \hspace{1.5cm} = & \textit{Automatic Voltage Regulator} \\
	DAE      & \hspace{1.5cm} = & \textit{Differential-Algebraic Equations} \\
	DFIG     & \hspace{1.5cm} = & \textit{Doubly Fed Induction Generator} \\
	DFIG-WT  & \hspace{1.5cm} = & \textit{Doubly Fed Induction Generator Wind Turbine} \\
	DR       & \hspace{1.5cm} = & \textit{Damping Ratio} \\
	EBT      & \hspace{1.5cm} = & Energi Baru Terbarukan \\
	FRC      & \hspace{1.5cm} = & \textit{Full-Rated Converter} \\
	GSC      & \hspace{1.5cm} = & \textit{Grid Side Converter} \\
	IBR      & \hspace{1.5cm} = & \textit{Inverter-Based Resources} \\
	IEEE     & \hspace{1.5cm} = & \textit{Institute of Electrical and Electronics Engineers} \\
	IRC      & \hspace{1.5cm} = & \textit{Inner Rotor Current Controller} \\
	PLL      & \hspace{1.5cm} = & \textit{Phase-Locked Loop} \\
	PLN      & \hspace{1.5cm} = & Perusahaan Listrik Negara \\
	PLTB     & \hspace{1.5cm} = & Pembangkit Listrik Tenaga Bayu \\
	PMSG     & \hspace{1.5cm} = & \textit{Permanent Magnet Synchronous Generator} \\
	PMSG-WT  & \hspace{1.5cm} = & \textit{Permanent Magnet Synchronous Generator Wind Turbine} \\
	PQC      & \hspace{1.5cm} = & \textit{Active and Reactive Power Controller} \\
	PRF      & \hspace{1.5cm} = & \textit{Participation Factor} \\
	PSS      & \hspace{1.5cm} = & \textit{Power System Stabilizer} \\
	RSC      & \hspace{1.5cm} = & \textit{Rotor Side Converter} \\
	RUPTL    & \hspace{1.5cm} = & Rencana Usaha Penyediaan Tenaga Listrik \\
	SC       & \hspace{1.5cm} = & \textit{Speed Controller} \\
	SCIG     & \hspace{1.5cm} = & \textit{Squirrel Cage Induction Generator} \\
	SMIB     & \hspace{1.5cm} = & \textit{Single Machine Infinite Bus} \\
	SSS      & \hspace{1.5cm} = & \textit{Small Signal Stability} \\
	WECC     & \hspace{1.5cm} = & \textit{Western Electricity Coordinating Council} \\
	WTG      & \hspace{1.5cm} = & \textit{Wind Turbine Generator} \\
\end{tabular}
